% Clase
\documentclass[11pt,a4paper,spanish,twoside]{report}

% Órdenes auxiliares
% Español
\usepackage[spanish]{babel}
\usepackage[utf8]{inputenc}
\usepackage[T1]{fontenc}
\usepackage{lmodern}

% Imágenes
\usepackage[pdftex]{graphicx}
\usepackage{latexsym}
\usepackage{fancybox}

% Ruta para las imágenes
\graphicspath{{img/}}

% Rotaciones
\usepackage[twoside]{rotating}

% Multirow
\usepackage{multirow}

% Referencias
\usepackage[spanish]{varioref}
\usepackage[pdftex,colorlinks=true,linkcolor=black]{hyperref}

% Colores
\usepackage{color}
\usepackage{colortbl}

% Párrafos
\setlength{\parskip}{6pt}

% Code for creating empty pages
% No headers on empty pages before new chapter
\makeatletter
\def\cleardoublepage{\clearpage\if@twoside \ifodd\c@page\else
    \hbox{}
    \thispagestyle{empty}
    \newpage
    \if@twocolumn\hbox{}\newpage\fi\fi\fi}
\makeatother \clearpage{\pagestyle{empty}\cleardoublepage}

%% Entorno Listings para código fuente
\usepackage{listingsutf8}[2007/11/11]

\lstset{
  frame=Ltb, %forma del cuadro
  framerule=0pt, %ancho del borde
  aboveskip=0.5cm, %separación de los números de línea
  framexleftmargin=0.4cm, %margen externo izquierdo
  framesep=0pt,
  rulesep=.4pt,
  rulesepcolor=\color{black},
  % 
  stringstyle=\ttfamily,
  showstringspaces = false,
  basicstyle=\scriptsize,
  keywordstyle=\bfseries,
  % 
  numbers=left,
  numbersep=15pt,
  numberstyle=\tiny,
  numberfirstline= false,
  %
  inputencoding=utf8/latin1
}

% minimizar fragmentado de listados
\lstnewenvironment{listing}[1][]{
  \lstset{#1}\pagebreak[0]}{\pagebreak[0]
}


% Fancyhdr - Encabezados y pies de página
\usepackage{fancyhdr}
% Márgenes
\headsep=8mm
\footskip=14mm

% Fancy Header Style Options
\pagestyle{fancy}               % Sets fancy header and footer
\fancyfoot{}                    % Delete current footer settings

% Sin mayúsculas en la cabecera
\lhead{\nouppercase{\rightmark}}
\rhead{\nouppercase{\leftmark}}

% Capítulo
\renewcommand{\chaptermark}[1]{ % Lower Case Chapter marker style
  \markboth{\chaptername\ \thechapter.\ #1}{}} 

% Sección
\renewcommand{\sectionmark}[1]{ % Lower case Section marker style
  \markright{\thesection.\ #1}} 

% Página
\fancyhead[LE,RO]{\bfseries\thepage} % Page number (boldface) in left on even
                                     % pages and right on odd pages

% ------------------ Macro para encabezados y pies de página-------------------
%    Uso: \encabezado{pares(pag izquierda)}
% -----------------------------------------------------------------------------
\def\encabezado{
  \fancyhead[RE]{\bfseries\leftmark}     % In the right on even pages
  \fancyhead[LO]{\bfseries\rightmark}  % In the left on odd pages
  \renewcommand{\headrulewidth}{0.5pt} % Width of head rule
}
% -----------------------------------------------------------------------------

% ------------------ Macro para insertar una imagen ---------------------------
%    Uso: \imagen{nombreFichero}{Ancho(cm))}{Etiqueta}{Identificador}
% -----------------------------------------------------------------------------
\usepackage{float}
\usepackage{ifthen}
\def\imagen#1#2#3#4{
  \begin{figure}[H]
    \begin{center}
      \ifthenelse{\equal{#2}{}}
      {\includegraphics{#1}}{\resizebox{#2cm}{!}{\includegraphics{#1}}}
      \ifthenelse{\equal{#3}{}}{}{\caption{#3}}
      \label{#4}
    \end{center}
  \end{figure}
}
% -----------------------------------------------------------------------------


% ------------------ Macro para la portada ------------------------------------
%    Uso: \portada{asignatura}{titulo}{subtítulo}{autor}{fecha}
% -----------------------------------------------------------------------------
\def\portada#1#2#3#4#5{
  \thispagestyle{empty}
  \vspace*{-3.3cm}

  \begin{minipage}[t]{14cm}
    \begin{center}
      \includegraphics[scale=0.25]{logoesi}
  
      \vspace*{1.5cm}
      {\Large \textbf{Universidad de Castilla-La Mancha\\ 
          Escuela Superior de Informática}\\}
    
      \vspace*{1.2cm}
      {\huge \textbf{#1}\\}
    
      \vspace*{1.5cm}
      {\huge #2}\\{\Large #3\\}
    
      \vspace*{1.5cm}
      {\large #4\\}
      \vspace*{1.4cm}
      \large{#5}
    \end{center}
  \end{minipage}

  \newpage
%   \vspace*{1cm}
%   \thispagestyle{empty} 
%   \newpage
}
% -----------------------------------------------------------------------------


% ------------------ Macro para la licencia -----------------------------------
%    Uso: \portada{asignatura}{titulo}{subtítulo}{autor}{fecha}
% -----------------------------------------------------------------------------
\def\licencia#1{
  \thispagestyle{empty}  % Suprime la numeración de esta página
  \vspace*{16.5cm}
  \begin{small}
    \copyright~ #1. Se permite la copia, distribución y/o 
    modificación de este documento bajo los términos de la licencia de
    documentación libre GNU, versión 1.1 o cualquier versión posterior publicada
    por la {\em Free Software Foundation}, sin secciones invariantes. Puede
    consultar esta licencia en http://www.gnu.org. \\[0.2cm]
    Este documento fue compuesto con \LaTeX{}. 
  \end{small}
  \newpage
%   \thispagestyle{empty}
%   \vspace*{0cm}
%   \newpage
}
% -----------------------------------------------------------------------------




% Encabezado y pie de página
\encabezado

\setcounter{secnumdepth}{3}
\setcounter{tocdepth}{3}

\begin{document}

% Silabación extra
\hyphenation{
a-sig-na-tu-ras
au-to-ma-ti-za-rá
ca-tá-lo-go
ca-rre-ra
diag-nos-tico
in-fe-rior
man-te-ni-mien-to
per-so-nal
pro-por-cio-na-rá
pu-bli-ca-da
re-qui-si-tos
res-pecto
u-su-a-rios
vi-lla-rre-al
}


% Portada
\portada{Sistemas para la colaboración}
{Trabajo teórico:}{Sistemas colaborativos para la gestión de proyectos}
{Sergio de la Rubia García-Carpintero\\Alicia Martín-Benito Escalona}
{28 de Abril de 2011}

% Licencia
\licencia{Sergio de la Rubia García-Carpintero, Alicia Martín-Benito Escalona}

% Índices
\tableofcontents
% \listoffigures
% \listoftables

%% INICIO DEL DOCUMENTO %%%%%%%%%%%%%%%%%%%%%%%%%%%%%%%%%%%%%%%%%%%%%%%%%

\chapter{Introducción}
La gestión de proyectos es la disciplina que guía e integra los procesos de
planificar, captar, dinamizar, organizar talentos y administrar recursos, con
el fin de culminar todo el trabajo requerido para desarrollar un proyecto y
cumplir con el alcance, dentro de unos límites de tiempo y costo definidos.
Todo lo cual requiere liderar los talentos, evaluar y regular continuamente
las acciones necesarias y suficientes \cite{GPR}.

Desde los tiempos de los primeros ordenadores comerciales, las empresas
buscaban sistemas informáticos que les ayudaran a planificar y gestionar sus
proyectos. Estos sistemas llegaron a ser muy potentes pero eran difíciles de
usar y poco amigables. Además, eran sistemas no colaborativos y sus datos
solían estar solo disponibles en los terminales de la propia empresa.

Con el paso del tiempo, dichos sistemas software han ido evolucionando y han
ido incorporando nuevas funcionalidades: manejo y control de presupuestos,
asignación de recursos, manejo de la calidad, software para la colaboración,
software para la comunicación, etc.

Actualmente, existe una gama muy amplia de aplicaciones (privativas y no
privativas) para la gestión de proyectos. Estas aplicaciones suelen
clasificarse atendiendo a las funcionalidades que implementan. Las principales
funcionalidades implementadas por estas son las siguientes \cite{WPM}:
\begin{itemize}
	\item Software colaborativo.
	\item Sistema de seguimiento de incidentes.
	\item Planificación.
	\item Gestión de la cartera de proyectos.
	\item Gestión de recursos.
	\item Gestión de documentos.
\end{itemize}
Además, estas aplicaciones pueden o no estar basadas en web.

El objetivo de este trabajo es analizar de entre todas las aplicaciones
disponibles, aquellas que tienen la funcionalidad de software colaborativo.
Es decir, aquellas aplicaciones que permiten trabajar en un mismo proyecto y
de forma concurrente a varios usuarios que se encuentran en diversas
estaciones de trabajo, conectados a través de una red.

Para ello, en primer lugar se explicará en qué consiste y cómo está
implementada cada una de las funcionalidades, en alguno de los programas más
utilizados actualmente. Y después, para un mejor entendimiento de cómo todas
estas funcionalidades conviven en un mismo software, se estudiará el caso
particular de la aplicación \emph{Redmine}.

\chapter{Sistema de seguimiento de incidentes}
\section{En qué consiste}
Es un paquete software que mantiene y adminstra una lista de incidentes que
han sido reportados. También mantiene una base de datos con información sobre
problemas comunes y su solución.

\section{Ejemplo. \emph{Trac}}
  \subsection{Información general}
    \begin{itemize}
        \item\textbf{Desarrollado por:} \emph{Edgewall Software}
        \item\textbf{Sistema operativo:} Multiplataforma.
	\item\textbf{Lenguaje de implementación:} \emph{Python}.
        \item\textbf{Software colaborativo:} Si.
        \item\textbf{Sistema de seguimiento de incidentes:} Si.
	\item\textbf{Planificación:} No.
	\item\textbf{Gestión de la cartera de proyectos:} No.
	\item\textbf{Gestión de recursos:} No.
	\item\textbf{Gestión de documentos:} No.
	\item\textbf{Basado en web:} Si.
	\item\textbf{Licencia:} \emph{BSD modificada}.
	\item\textbf{URL:} \url{http://trac.edgewall.org/}
        \end{itemize}
	\subsection{Características}
	\subsection{Cómo lleva a cabo el seguimiento de incidentes}

\chapter{Planificación}
\section{En qué consiste}
La planificación es un proceso en el que se establece el esfuerzo necesario
a realizar para cumplir con unos objetivos en un tiempo determinado. Este
proceso permite además, refinar los objetivos que dieron origen al proyecto.

Si bien la planificación define las acciones a seguir durante un periodo de
tiempo, durante la ejecución de las mismas, puede surgir la necesidad de
realizar alguna modificación respecto a lo definido originalmente. En ese
caso, se realizará una nueva planificación partiendo de ese mismo instante y
teniendo en cuenta las acciones que aún no fueron realizadas.

\section{Ejemplo. \emph{Jira}}
\imagen{jira-logo.png}{5}{}{}
Es una aplicación basada en web que se utiliza para el seguimiento de errores,
de incidentes y para la gestión operativa de proyectos. También se utiliza en
áreas no técnicas para la administración de tareas. La herramienta fue
desarrollada por la empresa australiana \emph{Atlassian}. \cite{WJI}

  \subsection{Información general}
    \begin{itemize}
		\item \textbf{Desarrolado por:} \emph{Atlassian}
		\item \textbf{Sistema operativo:} Multiplataforma.
		\item \textbf{Lenguaje de implementación:} \emph{Java}.
    \item \textbf{Software colaborativo:} Si.
    \item \textbf{Sistema de seguimiento de incidentes:} Si.
		\item \textbf{Planificación:} Si.
		\item \textbf{Gestión de la cartera de proyectos:} No.
		\item \textbf{Gestión de recursos:} No.
		\item \textbf{Gestión de documentos:} No.
		\item \textbf{Basado en web:} Si.
		\item \textbf{Licencia:} Propietaria. Gratuita para uso privado.
		\item \textbf{URL:} \url{http://www.atlassian.com/software/jira}
		\end{itemize}

	\subsection{Características}
	\emph{Jira} está basado en \emph{Java EE}, que funciona en varias bases de
	datos y	sistemas operativos. La herramienta dispone también de paneles de
	control	adaptables,	filtros de búsqueda, estadísticas, RSS y función de
	correo electrónico.

	La flexible arquitectura de \emph{Jira} permite al usuario crear ampliaciones
	específicas que pueden incluirse en la \emph{Jira extension library}.

	A pesar de que \emph{Jira} es un producto comercial, que cuenta con clientes
	tan	importantes como: \emph{BMW}, \emph{Yahoo}, \emph{Adobe}, \emph{Boeing},
	\emph{Electronic Arts}, etc., se dan licencias gratuitas para proyectos
	\emph{Open-Source}, instituciones sin ánimo de lucro, organizaciones
	caritativas y personas individuales.

	\subsection{Cómo lleva a cabo la planificación}

\chapter{Gestión de la cartera de proyectos}
\section{En qué consiste}
Es un término utilizado por los jefes de proyecto y por las organizaciones
gestoras de proyectos, para describir los métodos para el análisis y la
gestión colectiva de un grupo de proyectos propuestos o en curso.
El objetivo fundamental es determinar cuál es la combinación de los proyectos
y la secuencia de realización que estos deben seguir para lograr los mejores
resultados según los objetivos de la organización.

\section{Ejemplo. \emph{Launchpad}}
\imagen{launchpad-logo.png}{6}{}{}
Es una aplicación web que permite desarrollar y mantener software, en
particular software libre. Está desarrollada y mantenida por \emph{Canonical
Ltd.} En 2009 \emph{Launchpad} pasó a ser completamente libre, bajo la
versión 3 de la licencia \emph{AGPL} (\emph{Affero General Public License}).
\cite{WLA}

  \subsection{Información general}
    \begin{itemize}
		\item \textbf{Desarrolado por:} \emph{Canonical Ltd.}
		\item \textbf{Sistema operativo:} Multiplataforma.
		\item \textbf{Lenguaje de implementación:} \emph{Python}.
    \item \textbf{Software colaborativo:} Si.
    \item \textbf{Sistema de seguimiento de incidentes:} Si.
		\item \textbf{Planificación:} No.
		\item \textbf{Gestión de la cartera de proyectos:} Si.
		\item \textbf{Gestión de recursos:} No.
		\item \textbf{Gestión de documentos:} No.
		\item \textbf{Basado en web:} Si.
		\item \textbf{Licencia:} \emph{Affero GPL}.
		\item \textbf{URL:} \url{https://launchpad.net/launchpad-project}
		\end{itemize}

	\subsection{Características}
	Consta de varias partes:
	\begin{itemize}
	\item \textbf{Code:} un sitio de alojamiento de código fuente que utiliza el
	sistema de control de versiones \emph{Bazaar}.
	\item \textbf{Bugs:} un sistema de seguimiento de errores para informar
	sobre \emph{bugs} en diferentes distribuciones y productos.
	\item \textbf{Blueprints:} un sistema de seguimiento para especificaciones y
	nuevas características.
	\item \textbf{Translations:} un sitio para traducir aplicaciones a múltiples
	idiomas.
	\item \textbf{Answers:} un sitio de ayuda para la comunidad.
	\item \textbf{Soyuz:} una herramienta para llevar una pequeña parte del
	mantenimiento de las distribuciones. Abarca el sistema de construcción, el
	mantenimiento de paquetes y la publicación de archivos.
	\end{itemize}

	\emph{Launchpad} es usada primordialmente para el desarrollo de
	\emph{Ubuntu} y sus	derivados oficiales, aunque también contempla otras
	distribuciones y proyectos independientes.

	\subsection{Cómo lleva a cabo la gestión de la cartera de proyectos}

\chapter{Gestión de recursos}
\section{En qué consiste}
\section{Ejemplo. \emph{MindGenius}}
  \subsection{Información general}
    \begin{description}
        \item\textbf{Desarrollado por:} \emph{MindGenius Ltd.}
	\item\textbf{Sistema operativo:} \emph{Microsoft Windows}.
	\item\textbf{Lenguaje de implementación:} .
        \item\textbf{Software colaborativo:} Si. 
        \item\textbf{Sistema de seguimiento de incidentes:} No. 
	\item\textbf{Planificación:} Si.
	\item\textbf{Gestión de la cartera de proyectos:} No.
	\item\textbf{Gestión de recursos:} Si.
	\item\textbf{Gestión de documentos:} No.
	\item\textbf{Basado en web:} No.
	\item\textbf{Licencia:} \emph{Comercial}.
	\item\textbf{URL:} \url{http://www.mindgenius.com/}.
	\end{description}
	\subsection{Características}
	\subsection{Cómo lleva a cabo la gestión de recursos}

\chapter{Gestión de documentos}
\section{En qué consiste}
\section{Ejemplo. \emph{5pm}}
  \subsection{Información general}
    \begin{description}
        \item\textbf{Desarrollado por:} \emph{Quatre Group LLC.}
	\item\textbf{Sistema operativo:} Multiplataforma. 
	\item\textbf{Lenguaje de implementación:} \emph{AJAX} y \emph{Adobe Flash}.
        \item\textbf{Software colaborativo:} Si. 
        \item\textbf{Sistema de seguimiento de incidentes:} No. 
	\item\textbf{Planificación:} Si.
	\item\textbf{Gestión de la cartera de proyectos:} No.
	\item\textbf{Gestión de recursos:} No.
	\item\textbf{Gestión de documentos:} Si.
        \item\textbf{Basado en web:} Si.
	\item\textbf{Licencia:} \emph{SaaS}.
	\item\textbf{URL:} \url{www.5pmweb.com}
	\end{description}
	\subsection{Características}
	\subsection{Cómo lleva a cabo la gestión de documentos}

\chapter{Caso de estudio. \emph{Redmine}}
\imagen{redmine-logo.png}{8}{}{}
\section{Información general}
	\begin{itemize}
		\item \textbf{Desarrolado por:} \emph{Jean-Philippe Lang y otros}
		\item \textbf{Sistema operativo:} Multiplataforma.
		\item \textbf{Lenguaje de implementación:} \emph{Ruby on Rails}.
    \item \textbf{Software colaborativo:} Si.
    \item \textbf{Sistema de seguimiento de incidentes:} Si.
		\item \textbf{Planificación:} Si.
		\item \textbf{Gestión de la cartera de proyectos:} Si.
		\item \textbf{Gestión de recursos:} No.
		\item \textbf{Gestión de documentos:} Si.
		\item \textbf{Basado en web:} Si.
		\item \textbf{Licencia:} \emph{GPL}v2.
		\item \textbf{URL:} \url{http://www.redmine.org}
	\end{itemize}
\section{Características}
\section{Funcionamiento}
	\subsection{Sistema de seguimiento de incidentes}
	\subsection{Planificación}
	\subsection{Gestión de la cartera de proyectos}
	\subsection{Gestión de documentos}

\begin{thebibliography}{99}
\bibitem{GPR} Definición de \emph{Gestión de proyectos}:
\\ \url{http://es.wikipedia.org/wiki/Gestion_de_proyectos}
\bibitem{WPM} Comparación del software disponible para la gestión de proyectos:
\\ \url{http://en.wikipedia.org/wiki/List_of_project_management_software}
\bibitem{WTR} Descripción de Trac:
  \url{http://es.wikipedia.org/wiki/Trac}
\bibitem{TRA} Página web de Trac:
  \url{http://trac.edgewall.org/}
\bibitem{WJI} Descripción de Jira:
  \url{http://es.wikipedia.org/wiki/JIRA}
\bibitem{JIR} Página web de Jira:
  \url{http://www.atlassian.com/software/jira}
\bibitem{WLA} Descripción de Launchpad:\\
  \url{http://es.wikipedia.org/wiki/Launchpad}
\bibitem{LAU} Página web de Launchpad:\\
  \url{http://launchpad.net/launchpad-project}
\bibitem{WMG} Descripción de MindGenius:\\
  \url{http://en.wikipedia.org/wiki/MindGenius}
\bibitem{MGE} Página web de MindGenius:
  \url{http://www.mindgenius.com}
\bibitem{W5P} Descripción de 5pm:
  \url{http://en.wikipedia.org/wiki/5pm}
\bibitem{5PM} Página web de 5pm:
  \url{http://www.5pmweb.com}
\bibitem{WRM} Descripción de Redmine:
  \url{http://es.wikipedia.org/wiki/Redmine}
\bibitem{RMI} Página web de Redmine:
  \url{http://www.redmine.org}
\end{thebibliography}

\end{document}
